% ============================================================
\chapter{Connexions et D\'eriv\'ees Covariantes}
\label{ch:connexions}
% ============================================================

Comment d\'eriver un champ de vecteurs sur une sph\`ere~? Sur $\R^n$, on
d\'erive composante par composante. Mais sur une vari\'et\'e courbe, les
espaces tangents changent de point en point, et comparer deux vecteurs en
des points diff\'erents n'a aucun sens a priori. Il faut un ingr\'edient
suppl\'ementaire~: une \emph{connexion}, qui prescrit comment
\og transporter\fg{} un vecteur le long d'une courbe. C'est Tullio
Levi-Civita qui, en 1917, a formalis\'e cette id\'ee dans le contexte
riemannien, en montrant qu'il existe une unique connexion compatible avec
la m\'etrique et sans torsion.

% ------------------------------------------------------------
\section{Motivation~: la d\'eriv\'ee directionnelle ne suffit pas}
% ------------------------------------------------------------

Sur $\R^n$, la d\'eriv\'ee d'un champ de vecteurs $Y$ dans la direction $X$ est
simplement $(D_X Y)^i = X^j \frac{\partial Y^i}{\partial x^j}$. Sur une vari\'et\'e
g\'en\'erale, cette op\'eration n'est pas intrins\`eque~: elle d\'epend du choix de
coordonn\'ees. La notion de \emph{connexion} fournit le cadre correct pour d\'eriver
des sections de fibr\'es vectoriels.

\section{Connexions sur les fibr\'es vectoriels}

\begin{definition}[Connexion lin\'eaire]
Soit $E \to M$ un fibr\'e vectoriel lisse. Une \emph{connexion} (ou \emph{d\'eriv\'ee
covariante}) sur $E$ est une application $\R$-bilin\'eaire~:
\[
    \nabla : \mathfrak{X}(M) \times \Gamma(E) \to \Gamma(E),
    \quad (X, s) \mapsto \nabla_X s,
\]
satisfaisant les axiomes suivants~:
\begin{enumerate}
    \item \textbf{$C^\infty(M)$-lin\'earit\'e en $X$~:} $\nabla_{fX} s = f\, \nabla_X s$
          pour tout $f \in C^\infty(M)$.
    \item \textbf{R\`egle de Leibniz~:} $\nabla_X(fs) = (Xf)\, s + f\, \nabla_X s$
          pour tout $f \in C^\infty(M)$.
\end{enumerate}
\end{definition}

\begin{remarque}
L'axiome (1) montre que $\nabla_X s$ en un point $p$ ne d\'epend que de $X_p$
(et non du champ $X$ tout entier). L'axiome (2) montre que $\nabla$ n'est
\emph{pas} $C^\infty(M)$-lin\'eaire en $s$~: c'est un op\'erateur diff\'erentiel
d'ordre $1$.
\end{remarque}

\begin{proposition}[Localit\'e]
Si $s_1 = s_2$ sur un ouvert $U \subset M$, alors $\nabla_X s_1 = \nabla_X s_2$
sur $U$. La connexion est donc un op\'erateur local.
\end{proposition}

\section{Expression en coordonn\'ees~: symboles de Christoffel}

Soit $(U, x^1, \ldots, x^n)$ une carte locale et $(e_1, \ldots, e_r)$ un rep\`ere
local de $E$ sur $U$. On d\'efinit les \emph{coefficients de connexion}
$\omega^k_{ij}$ par~:
\[
    \nabla_{\partial_i} e_j = \omega^k_{ij}\, e_k.
\]

Pour le fibr\'e tangent $E = TM$, ces coefficients sont not\'es $\Gamma^k_{ij}$
et appel\'es \emph{symboles de Christoffel}~:
\[
    \nabla_{\partial_i} \partial_j = \Gamma^k_{ij}\, \partial_k.
\]

\begin{proposition}[Formule en coordonn\'ees]
Si $Y = Y^j \partial_j$ et $X = X^i \partial_i$, alors~:
\[
    \nabla_X Y = X^i \left(\frac{\partial Y^k}{\partial x^i} + \Gamma^k_{ij} Y^j\right) \partial_k.
\]
On note la composante~: $(\nabla_X Y)^k = X^i (\partial_i Y^k + \Gamma^k_{ij} Y^j)
= X^i \nabla_i Y^k$.
\end{proposition}

\begin{proof}
Par la r\`egle de Leibniz~:
\begin{align*}
    \nabla_X Y &= \nabla_{X^i \partial_i}(Y^j \partial_j)
    = X^i \nabla_{\partial_i}(Y^j \partial_j) \\
    &= X^i \left[(\partial_i Y^j)\, \partial_j + Y^j \nabla_{\partial_i}\partial_j\right]
    = X^i \left[(\partial_i Y^j)\, \partial_j + Y^j \Gamma^k_{ij}\, \partial_k\right].
\end{align*}
En renommant l'indice muet~: $(\nabla_X Y)^k = X^i(\partial_i Y^k + \Gamma^k_{ij} Y^j)$.
\end{proof}

\section{Changement de cartes}

\begin{proposition}[Transformation des symboles de Christoffel]
Sous un changement de coordonn\'ees $x^i \mapsto \bar{x}^i$, les symboles de
Christoffel se transforment selon~:
\[
    \bar{\Gamma}^k_{ij} = \frac{\partial \bar{x}^k}{\partial x^l}
    \frac{\partial x^m}{\partial \bar{x}^i} \frac{\partial x^n}{\partial \bar{x}^j}
    \Gamma^l_{mn}
    + \frac{\partial \bar{x}^k}{\partial x^l}
    \frac{\partial^2 x^l}{\partial \bar{x}^i \partial \bar{x}^j}.
\]
Le terme inhomog\`ene montre que les $\Gamma^k_{ij}$ ne forment \emph{pas} un tenseur.
\end{proposition}

\begin{intuition}[title=\textbf{Les symboles de Christoffel ne sont pas des tenseurs}]
La pr\'esence du terme $\frac{\partial^2 x^l}{\partial \bar{x}^i \partial \bar{x}^j}$
refl\`ete le fait que la connexion encode une information suppl\'ementaire par rapport
\`a la structure de vari\'et\'e seule. C'est pr\'ecis\'ement cette <<~correction~>>
qui permet de d\'eriver covariamment.
\end{intuition}

\section{Transport parall\`ele}

\begin{definition}[Champ parall\`ele le long d'une courbe]
Soit $\gamma : [a,b] \to M$ une courbe lisse et $V(t)$ un champ de vecteurs le
long de $\gamma$ (i.e.\ $V(t) \in T_{\gamma(t)}M$). On dit que $V$ est
\emph{parall\`ele} le long de $\gamma$ si~:
\[
    \frac{D V}{dt} := \nabla_{\dot{\gamma}(t)} V = 0.
\]
\end{definition}

En coordonn\'ees, la condition de parall\'elisme s'\'ecrit~:
\[
    \frac{dV^k}{dt} + \Gamma^k_{ij}(\gamma(t))\, \dot{\gamma}^i(t)\, V^j(t) = 0,
    \quad k = 1, \ldots, n.
\]
C'est un syst\`eme d'\'equations diff\'erentielles lin\'eaires du premier ordre.

\begin{theoreme}[Existence et unicit\'e du transport parall\`ele]
Soit $\gamma : [a,b] \to M$ une courbe lisse et $v_0 \in T_{\gamma(a)}M$.
Il existe un unique champ parall\`ele $V(t)$ le long de $\gamma$ avec $V(a) = v_0$.
L'application~:
\[
    P_\gamma^{a \to b} : T_{\gamma(a)}M \to T_{\gamma(b)}M,
    \quad v_0 \mapsto V(b),
\]
est un isomorphisme lin\'eaire, appel\'e \emph{transport parall\`ele} le long de $\gamma$.
\end{theoreme}

\begin{proof}
L'existence et l'unicit\'e d\'ecoulent du th\'eor\`eme de Cauchy--Lipschitz
appliqu\'e au syst\`eme lin\'eaire ci-dessus. La lin\'earit\'e de $P_\gamma^{a \to b}$
r\'esulte de la lin\'earit\'e du syst\`eme. L'inversibilit\'e s'obtient en consid\'erant
le transport parall\`ele le long de $\gamma$ parcourue en sens inverse.
\end{proof}

\begin{exemple}[Transport parall\`ele sur $S^2$]
Consid\'erons le triangle sph\'erique form\'e par des arcs de grands cercles reliant
le p\^ole Nord $N = (0,0,1)$ au point $A = (1,0,0)$, puis $A$ \`a $B = (0,1,0)$,
puis $B$ \`a $N$. Le transport parall\`ele d'un vecteur le long de ce triangle
le fait tourner d'un angle $\pi/2$. Cet angle est \'egal \`a l'aire du triangle
sph\'erique, illustrant le th\'eor\`eme de Gauss--Bonnet.
\end{exemple}

\section{Torsion d'une connexion}

\begin{definition}[Tenseur de torsion]
La \emph{torsion} d'une connexion $\nabla$ sur $TM$ est le tenseur de type $(1,2)$~:
\[
    T(X, Y) = \nabla_X Y - \nabla_Y X - [X, Y].
\]
\end{definition}

\begin{proposition}
$T$ est bien un tenseur~: il est $C^\infty(M)$-bilin\'eaire et antisym\'etrique.
En coordonn\'ees~: $T^k_{ij} = \Gamma^k_{ij} - \Gamma^k_{ji}$.
\end{proposition}

\begin{proof}
V\'erifions la $C^\infty(M)$-lin\'earit\'e~:
\begin{align*}
    T(fX, Y) &= \nabla_{fX} Y - \nabla_Y(fX) - [fX, Y] \\
    &= f\nabla_X Y - f\nabla_Y X - (Yf)X - f[X,Y] + (Yf)X \\
    &= f\, T(X, Y).
\end{align*}
L'antisym\'etrie est imm\'ediate. En coordonn\'ees~:
$T(\partial_i, \partial_j) = \Gamma^k_{ij}\partial_k - \Gamma^k_{ji}\partial_k$
car $[\partial_i, \partial_j] = 0$.
\end{proof}

\begin{definition}[Connexion sans torsion]
Une connexion est dite \emph{sans torsion} (ou \emph{sym\'etrique}) si $T \equiv 0$,
c'est-\`a-dire~: $\nabla_X Y - \nabla_Y X = [X, Y]$ pour tous $X, Y$.
En coordonn\'ees~: $\Gamma^k_{ij} = \Gamma^k_{ji}$.
\end{definition}

\section{Holonomie}

\begin{definition}[Groupe d'holonomie]
Soit $(M, \nabla)$ une vari\'et\'e munie d'une connexion et $p \in M$. Le
\emph{groupe d'holonomie} de $\nabla$ en $p$ est~:
\[
    \operatorname{Hol}_p(\nabla) = \{ P_\gamma : T_pM \to T_pM :
    \gamma \text{ lacet lisse en } p \} \subset GL(T_pM).
\]
Le \emph{groupe d'holonomie restreint} $\operatorname{Hol}^0_p(\nabla)$ est le
sous-groupe correspondant aux lacets contractiles.
\end{definition}

\begin{theoreme}[Ambrose--Singer]
Le groupe d'holonomie restreint $\operatorname{Hol}^0_p(\nabla)$ est un sous-groupe
de Lie connexe de $GL(T_pM)$ dont l'alg\`ebre de Lie est engendr\'ee par les
\'el\'ements de la forme $P_\gamma^{-1} \circ R(X, Y) \circ P_\gamma$,
o\`u $\gamma$ parcourt les chemins de $p$ \`a un point $q$ et $X, Y \in T_qM$.
\end{theoreme}

\begin{exemple}[Holonomie de $S^n$]
Pour la sph\`ere $S^n$ ($n \geq 2$), $\operatorname{Hol}(S^n) = SO(n)$, car la
sph\`ere est un espace sym\'etrique irr\'eductible simplement connexe.
\end{exemple}

\section{D\'eriv\'ee covariante des tenseurs}

\begin{definition}[Extension aux tenseurs]
Soit $\nabla$ une connexion sur $TM$. On \'etend $\nabla$ \`a tous les tenseurs par
les r\`egles~:
\begin{enumerate}
    \item $\nabla_X f = Xf$ pour $f \in C^\infty(M)$.
    \item $\nabla_X(S \otimes T) = (\nabla_X S) \otimes T + S \otimes (\nabla_X T)$
          (r\`egle de Leibniz).
    \item $\nabla$ commute avec les contractions.
    \item Sur les $1$-formes~: $(\nabla_X \omega)(Y) = X(\omega(Y)) - \omega(\nabla_X Y)$.
\end{enumerate}
\end{definition}

\begin{proposition}[D\'eriv\'ee covariante de la m\'etrique]
Pour un tenseur $(0,2)$ sym\'etrique $g$~:
\[
    (\nabla_X g)(Y, Z) = X(g(Y,Z)) - g(\nabla_X Y, Z) - g(Y, \nabla_X Z).
\]
On dit que $\nabla$ est \emph{compatible avec $g$} si $\nabla g = 0$, i.e.\
$(\nabla_X g)(Y,Z) = 0$ pour tous $X, Y, Z$.
\end{proposition}

\begin{formules}[title=\textbf{D\'eriv\'ee covariante en coordonn\'ees}]
Pour un tenseur de type $(r, s)$ avec composantes $T^{i_1 \cdots i_r}_{j_1 \cdots j_s}$~:
\[
    \nabla_k T^{i_1 \cdots i_r}_{j_1 \cdots j_s}
    = \partial_k T^{i_1 \cdots i_r}_{j_1 \cdots j_s}
    + \sum_{m=1}^r \Gamma^{i_m}_{k\ell}\, T^{i_1 \cdots \ell \cdots i_r}_{j_1 \cdots j_s}
    - \sum_{m=1}^s \Gamma^\ell_{k j_m}\, T^{i_1 \cdots i_r}_{j_1 \cdots \ell \cdots j_s}.
\]
\end{formules}

\section{Compatibilit\'e avec la m\'etrique et cons\'equences}

\begin{proposition}[Caract\'erisation de la compatibilit\'e]
Les conditions suivantes sont \'equivalentes~:
\begin{enumerate}
    \item $\nabla g = 0$.
    \item $X \inner{Y}{Z} = \inner{\nabla_X Y}{Z} + \inner{Y}{\nabla_X Z}$
          pour tous $X, Y, Z \in \mathfrak{X}(M)$.
    \item Le transport parall\`ele pr\'eserve le produit scalaire~: pour toute courbe
          $\gamma$ et tous $u, v \in T_{\gamma(a)}M$~:
          $\inner{P_\gamma u}{P_\gamma v}_{\gamma(b)} = \inner{u}{v}_{\gamma(a)}$.
    \item En coordonn\'ees~: $\partial_k g_{ij} = \Gamma^l_{ki} g_{lj} + \Gamma^l_{kj} g_{il}$.
\end{enumerate}
\end{proposition}

\begin{proof}
$(1) \Leftrightarrow (2)$~: C'est la d\'efinition de $\nabla g = 0$.

$(2) \Rightarrow (3)$~: Si $V(t)$ et $W(t)$ sont parall\`eles le long de $\gamma$,
alors $\frac{d}{dt}\inner{V}{W} = \inner{\nabla_{\dot\gamma}V}{W}
+ \inner{V}{\nabla_{\dot\gamma}W} = 0$, donc $\inner{V}{W}$ est constant.

$(3) \Rightarrow (2)$~: Tout champ de vecteurs le long de $\gamma$ s'\'ecrit
comme combinaison lin\'eaire de champs parall\`eles, et la condition s'en d\'eduit.
\end{proof}

\section{Espace affine des connexions}

\begin{proposition}
L'espace des connexions sur $TM$ est un espace affine model\'e sur
$\Gamma(T^*M \otimes T^*M \otimes TM)$~: si $\nabla$ et $\tilde{\nabla}$ sont
deux connexions, alors $A(X,Y) = \tilde{\nabla}_X Y - \nabla_X Y$ est un tenseur
de type $(1,2)$.
\end{proposition}

\begin{proof}
On v\'erifie la $C^\infty(M)$-bilin\'earit\'e~:
$A(fX, Y) = \tilde{\nabla}_{fX}Y - \nabla_{fX}Y = f(\tilde{\nabla}_X Y - \nabla_X Y) = fA(X,Y)$,
et $A(X, fY) = fA(X,Y) + (Xf)Y - (Xf)Y = fA(X,Y)$.
\end{proof}

\section{Exercices}

\begin{exercice}
Montrer que le transport parall\`ele le long d'un m\'eridien de $S^2$ (de $\theta_0$
\`a $\theta_1$, $\phi$ fix\'e) envoie $\partial_\phi$ sur $\partial_\phi$
(normalis\'e correctement).
\end{exercice}

\begin{exercice}
Calculer les symboles de Christoffel de la m\'etrique $g = dr^2 + r^2 d\theta^2$
en coordonn\'ees polaires sur $\R^2$. V\'erifier que la torsion est nulle.
\end{exercice}

\begin{exercice}
Soit $\nabla$ une connexion compatible avec la m\'etrique et $V(t)$ un champ
parall\`ele le long de $\gamma$. Montrer que $\norm{V(t)}$ est constant.
\end{exercice}

\begin{exercice}
Montrer que si deux connexions $\nabla$ et $\tilde{\nabla}$ sont toutes deux sans
torsion et compatibles avec la m\^eme m\'etrique $g$, alors $\nabla = \tilde{\nabla}$.
(C'est l'unicit\'e de la connexion de Levi-Civita, d\'emontr\'ee au chapitre suivant.)
\end{exercice}

\begin{exercice}
Soit $G$ un groupe de Lie avec m\'etrique invariante \`a gauche. Montrer que la
connexion d\'efinie par $\nabla_X Y = \frac{1}{2}[X, Y]$ pour $X, Y$ invariants
\`a gauche est sans torsion. Quand est-elle compatible avec la m\'etrique~?
\end{exercice}

\begin{exercice}
Calculer le groupe d'holonomie du tore plat $\mathbb{T}^2$. Comparer avec celui
de $S^2$.
\end{exercice}
